\documentclass[12pt]{article}
\usepackage[utf8]{inputenc}
\usepackage[T1]{fontenc}
\usepackage{hyperref}
\title{Transalpine Gaulish Language Resources for K-12 Curriculum Development}
\author{Compiled by Claude Sonnet 5.5 (claude-sonnet-5-5) for Andr\'as Kornai}
\date{}
\begin{document}
\maketitle

\section{Introduction}
Transalpine Gaulish (ISO 639-3 xtg) is an extinct Celtic language of ancient Gaul. It has no native speakers, so any K-12 use would be heritage and revitalization work built on reconstruction. The searches found good coverage of the neighbouring Cisalpine Celtic material but no standard Transalpine Gaulish corpus or dictionary. Cisalpine Gaulish is closer to Transalpine Gaulish than Lepontic is. The relationship between Lepontic and Gaulish is debated, and LexLep treats them as separate languages.

\section{Primary Educational Resources}
No school programmes or classroom materials were found.

\section{Linguistic and Reference Materials}
\begin{itemize}
\item Lexicon Leponticum (LexLep), University of Vienna: a digital edition of Cisalpine Celtic inscriptions with an etymological dictionary. It covers Lepontic and Cisalpine Gaulish, not Transalpine Gaulish. \url{https://lexlep.univie.ac.at/}
\item LexLep Dictionary of Attested Forms. \url{https://lexlep.univie.ac.at/wiki/Dictionary_of_Attested_Forms}
\item Digital Classicist wiki, Lexicon Leponticum. \url{https://wiki.digitalclassicist.org/Lexicon_Leponticum}
\item Stifter, ``Cisalpine Celtic'' (Palaeohispanica, 2020), an overview of about 430 short texts. \url{https://mural.maynoothuniversity.ie/id/eprint/14948/1/DS-Cisalpine-Celtic-2020.pdf}
\item Lexicon Leponticum (CLL 2024) PDF. \url{https://mural.maynoothuniversity.ie/id/eprint/19716/1/LexiconLeponticumCLL2024.pdf}
\end{itemize}

\section{Religious and Community Resources}
None found.

\section{Digital and Audio Resources}
A community wiki page of Gaulish learning resources mentions \textit{Rogouson Gallicon Rissouom} by Gerard Poitrenaud (2023), a companion to the reconstructed-Gaulish grammar \textit{Gall\^ica Iextis}. This is a secondary pointer from a fan-maintained page and has not been checked. \url{https://r.datuan.dev/r/Gaulish/wiki/index/resources}

\section{Curriculum Development Context}
Transalpine Gaulish is attested only in fragmentary inscriptions, so any curriculum would be reconstructive and should be presented that way. Suitable content includes Celtic history, comparative linguistics, short attested texts and place-name studies. Any Gaulish-language instruction would rest on reconstructed grammar from hobbyist and scholarly sources of varying reliability. A Transalpine corpus (for example inscription collections) should be sought separately.
\end{document}